\documentclass[tikz,border=5pt]{standalone}
\usepackage[UTF8]{ctex}
\usepackage{amsmath,amssymb}
\usetikzlibrary{arrows.meta,calc,angles,quotes,patterns,decorations.pathmorphing}
\definecolor{field}{HTML}{176B9B}
\definecolor{force}{HTML}{B64B32}
\tikzset{>=Latex,line width=.8pt,every node/.style={font=\small},field/.style={color=field},force/.style={color=force},body/.style={draw=black,fill=blue!6},ghost/.style={dashed,gray},labelbox/.style={fill=white,inner sep=1.5pt}}
\begin{document}
\begin{tikzpicture}
\fill[blue!3](0,0)ellipse(3.4 and 1.15);
\draw[ghost](3.4,0)arc(0:180:3.4 and 1.15);
\draw(-3.4,0)arc(180:360:3.4 and 1.15);
\shade[ball color=orange!65](0,0)circle(.42);
\shade[ball color=blue!45](2.6,-.74)circle(.13);
\draw[->](0,0)--(2.6,-.74)node[midway,below]{$R_P$};
\node[above]at(0,.55){CZ星};
\node[below right]at(2.6,-.74){行星};
\draw[->,field](3.9,0)--(5.8,.28)node[right]{观察者};
\draw[ghost](0,0)--(5.7,0);
\node[above]at(4.7,.5){视线与轨道面夹角 $\theta$};
\node[below]at(0,-1.65){空间轨道投影，星体尺寸不按比例};
\end{tikzpicture}
\begin{tikzpicture}
\shade[ball color=orange!65](0,0)circle(.65);
\node[below]at(0,-.8){$r_C,T_C$};
\draw[->,field](.8,.35)--(4.5,.35);
\draw[->,field](.8,-.35)--(4.5,-.35);
\draw[fill=blue!8](4.55,-.7)rectangle(4.65,.7);
\draw[<->](0,-1.4)--(4.6,-1.4)node[midway,below]{$l_C$};
\node[above]at(4.6,.85){单位接收面积};
\node[above]at(2.5,1.45){$S_C=\sigma T_C^4(r_C/l_C)^2$};
\end{tikzpicture}
\begin{tikzpicture}
\foreach\j in{0,1}{
\begin{scope}[yshift=-2.7*\j cm]
\draw[ghost](-1,-1)--(-1,.8);
\draw[body](-1,.35)circle(.22)node[above=10pt]{恒星};
\draw[body](-1,-.6)circle(.13)node[below=8pt]{行星};
\ifnum\j=0
\draw[->,field](-.73,.35)--(.1,.35);
\node[field]at(.8,.75){$+v_C\cos\theta$};
\draw[->,force](-1.2,-.6)--(-2,-.6);
\else
\draw[->,field](-1.27,.35)--(-2.1,.35);
\node[field]at(.8,.75){$-v_C\cos\theta$};
\draw[->,force](-.8,-.6)--(0,-.6);
\fi
\draw[->](2.2,0)--(3.7,0)node[right]{视线};
\end{scope}}
\node[below]at(1,-4.1){$\Delta\nu/\nu=2v_C\cos\theta/c$};
\node[above]at(1,1.1){两个径向速度极值时刻};
\end{tikzpicture}
\begin{tikzpicture}
\shade[ball color=blue!45](0,0)circle(1);
\foreach\y in{-.7,0,.7}\draw[->,field](-3,\y)--(-1.2,\y);
\node[below,align=center]at(-2,-1.1){入射功率\\$\pi r_P^2F$};
\foreach\a in{-70,-35,0,35,70}
\draw[->,force]({1.1*cos(\a)},{1.1*sin(\a)})--({1.65*cos(\a)},{1.65*sin(\a)});
\node[right,align=left]at(2,0){全表面辐射\\$4\pi r_P^2\sigma T_P^4$};
\draw[->](-.8,.6)--(-1.4,1.7)node[above]{$A_B$ 反射};
\end{tikzpicture}
\begin{tikzpicture}
\draw[body](0,0)circle(1.6);
\draw[ghost](0,0)circle(1.73);
\draw[->](-2.1,0)--(2.3,0)node[right]{$X$};
\draw[->](0,-1.9)--(0,2)node[above]{$Y$};
\draw[field,thick](-1.374,1.05)--(1.374,1.05);
\draw[body](-1.374,1.05)circle(.13);
\draw[body](1.374,1.05)circle(.13);
\draw[<->](-1.374,1.5)--(1.374,1.5)node[midway,above]{$v_PT_L$};
\draw[<->](.35,0)--(.35,1.05)node[midway,right,labelbox]{$b$};
\draw[ghost](0,0)--(1.374,1.05)node[pos=.73,below right,labelbox]{$r_C+r_P$};
\node[below]at(0,-2.4){接触判据作用于行星中心；尺寸为示意比例};
\end{tikzpicture}
\begin{tikzpicture}
\path[fill=blue!12](-1.766,.35)arc(180:360:1.766 and .22)
--(1.766,-.35)arc(0:-180:1.766 and .22)--cycle;
\draw(0,0)circle(1.8);
\draw[ghost](-1.766,.35)arc(180:0:1.766 and .22);
\draw[field](-1.766,.35)arc(180:360:1.766 and .22);
\draw[ghost](-1.766,-.35)arc(180:0:1.766 and .22);
\draw[field](-1.766,-.35)arc(180:360:1.766 and .22);
\draw[dash dot](0,-2.1)--(0,2.2)node[above]{轨道法线};
\draw[->](0,0)--(2.15,0);
\draw[->,field](0,0)--(1.766,.35);
\draw(1.3,0)arc(0:11.21:1.3);
\node[right]at(2,.4){$\theta_{\max}$};
\node[right,align=left]at(2.4,-.6){可凌日的视线带\\$|\theta|\le\theta_{\max}$\\$\eta=\sin\theta_{\max}$};
\fill(0,0)circle(1.5pt)node[below left]{$C$};
\end{tikzpicture}
\end{document}
